\cotefile{157-3}
\scanpage{66}
\mspage{15}

et soit \[(f_i)_{i\in I}:(X_i)_{i\in I}\longrightarrow(Y_i)_{i\in I}\] un \unsure{hom} de systèmes inductifs dans $\mathrm{Cat}$, indexés par un ens. ordonné filtrant, soient $X=\varinjlim X_i$, $Y=\varinjlim Y_i$, $f=\varinjlim f_i$, \deleted{$f:\varinjlim_i X_i\,\unsure{\to|=}$} $f:X\longrightarrow Y$. Si $f_i$ est dans $W$ pour tt $i$, alors $f\in W$.

On factorise cha\uline{que} $f_i$ en \[X_i\xrightarrow{\alpha_i}\overline X_i\xrightarrow{\overline f_i}Y_i,\qquad \alpha_i\in W_{\unsure{\omega|\infty}},\quad\overline f_i\in\mathrm{Prof}_{W_\omega}\]
% slanted
\margin{$W_\omega$ introduit dans la condition $4\,\mathrm b''$, page 14} \unsure{suivant} le système \unsure{décrit}. La construction étant fonctorielle, les $\overline X_i$ forment eux aussi un syst. inductif. De plus, comme le \deleted{\illegible{1}} foncteur \deleted{\illegible{1}} $S\longrightarrow\underline{\mathrm{Ch}}(S)$ qui donne naissance \deleted{\unsure{au}} à la factorisation canonique des morphismes, commute aux $\varinjlim$ filtrantes (immédiat), il s'ensuit que, si on pose \[\overline X=\varinjlim_i\overline X_{\unsure{i}},\] et qu'on regarde les flèches \[X\xrightarrow{\alpha}\overline X\xrightarrow{\overline f}Y,\qquad \alpha=\varinjlim\alpha_i,\quad\overline f=\varinjlim\overline f_i,\] avec $\overline f\alpha=f$, cette factorisation de $f$ est bien sa factorisation canonique. \inserted{Donc $\alpha\in W$, $\overline f\in\mathrm{Prof}_{W_\omega}$.} Ceci vu, en vertu des propriétés \deleted{\illegible{1}} $f\in W$, \deleted{\unsure{équivaut à dire que}} \inserted{i.e. $\overline f\in W$,} est vrai si les fibres de $\overline f$ \deleted{\unsure{sont}} \unsure{soient|sont} $W$-asphériques. Mais les fibres de $\overline f$ s'interprètent comme limites inductives filtrantes de fibres des $\overline f_i$. Lesquelles sont $W$-asphériques par \deleted{$4\,\unsure{\mathrm b}$} la condition $4\,\mathrm b''$ des fibrations. Donc l'axiome \textcircled{6} donne ce qu'il faut.

% slanted
\margin{et peut-être \textcircled{2} sous la forme plus \unsure{faible} \textcircled{\scriptsize 2bis}?} \uline{Conjecture} (\unsure{hardie}? \unsure{folle}?) Si $W\subset\mathrm{Fl}(\mathrm{Cat})$ satisfait les axiomes \textcircled{1} à \textcircled{6} (avec 4 sous sa forme la plus faible $4\,\mathrm b''$), alors $W=W_\infty$.

% archive sheet number: 66
\scanpage{67}
\mspage{16}

. Mais il est bien possible qu'il faille renforcer un peu $W(1)$ \inserted{(saturation)}, en exigeant que \deleted{\unsure{tout}} rétracté d'un $f\in W$ soit dans $W$.

\uline{\unsure{Remarques|Remarque}} Considérons, \inserted{pour $n\in\mathbb N$,} l'\unsure{ens} $W_n\subset\mathrm{Fl}(\mathrm{Cat})$ des ``$n$-équivalences \deleted{\unsure{d'homotopie}} \inserted{d'homotopie}'' de $\mathrm{Cat}$ : \[\begin{aligned}f\in W_n\iff{}&\pi_0(f)\text{ bijectif, et }\pi_i(f,x):\pi_i(X,x)\longrightarrow\pi_i(Y,fx)\\ &\text{bijectif pour }0\unsure{<|\leq}i\leq n.\end{aligned}\] Je croyais que ça \deleted{\unsure{fait}} satisfaisait à \deleted{\unsure{tout}} ce qu'on peut demander à un localiseur, sans avoir pris la peine de l'écrire avec \unsure{soin}. \deleted{\unsure{Bien sûr, ça sat'}} \inserted{Mais ça ne peut} \deleted{\unsure{pour plaire}} \inserted{satisfaire} à $W(4)$ \unsure{jusqu'à|puisqu'à} $W(4\,\mathrm b'')$, puisque on n'a pas $W_n\subset W_\infty$ !

$W(1)$ est trivial sur la définition (et on a sans bavure même la \uline{saturation forte} de $W_i$, i.e. si $f$ \deleted{\unsure{déjà}} \unsure{$\in\mathrm{Cat}$|t.q.} est iso dans $\mathrm{Hot}_{W_i}=(\mathrm{Cat})W_i^{-1}$, alors $f\in W_i$). Je crois me rappeler avoir vérifié $W(2)$, par la même méthode que pour $W_\infty$, en utilisant un critère cohomologique, à la Artin-Mazur, pour \deleted{$f\in$} exprimer $f\in W_i$, et en utilisant les suites spectrales de Leray pour $X\longrightarrow S$ et $Y\longrightarrow S$. La condition $W(3)$ est trivialement satisfaite puisque $W_n\supset W_\infty$ et qu'elle l'est pour $W_\infty$. Donc ça coince là où je m'y attendais le moins : l'axiome des fibrations. Prenons la forme faible : Soit $f:X\longrightarrow Y$ une $W_\omega$-fibration, \margin{$\omega>\infty$\\ $W_\omega\subset W_\infty\subset W_n$} et $f\in W_i$, les fibres\inserted{$X_s$} sont elles $W_i$-asphériques ? SPS $X$ et $Y$ 0-connexes, soit $x\in X_s$, la suite exacte d'homotopie donne \diagram{\begin{tikzcd}[ampersand replacement=\&,column sep=small] \pi_{i+1}(X)\arrow[r, "{\text{iso si }i\leq n-1}"']\&\pi_{i+1}(Y)\arrow[r]\&\pi_i(X_s)\arrow[r]\&\pi_i(X)\arrow[r,"{\mathrm{iso}}","{\text{si }i\leq n}"']\&\pi_i(Y) \end{tikzcd}} donc \[\pi_i(X_s)=1\qquad\text{si }0\leq i\leq n-1,\] mais pour $i=n$, on trouve

% archive sheet number: 67
\scanpage{68}
\mspage{17}

\[\pi_n(X_s)=\mathrm{Coker}\bigl(\pi_{n+1}(X)\longrightarrow\pi_{n+1}(Y)\bigr),\] il n'y a aucune raison que ce soit nul, comme ça devrait si $X_s\longrightarrow e$ était dans $W_n$.

Il est vrai que parfois \deleted{\unsure{on}} \unsure{renforce|renforcer} fortement la condition pour $W_n$, disons qu'on définit $W'_n\subset W_n$ par \[f\in W'_n\iff\left\{\begin{aligned}&\mathrm a)\ \pi_0(f)\text{ bijectif}\\&\mathrm b)\ \pi_i(f,x):\pi_i(X,x)\longrightarrow\pi_i(Y,fx)\text{ bij.}\\&\hspace{3em}\text{pour }0<i\leq n\text{ et }x\in X\\&\mathrm c)\ \pi_{n+1}(f,x):\pi_{\unsure{n+1}}(X,x)\to\pi_{n+1}(Y,fx)\ \text{\uline{surj}}.\end{aligned}\right.\] c'est à dire aussi \[f\in W'_n\iff\left\{\begin{aligned}&\text{les ``fibres homotopiques''}\ \inserted{Z}\ \text{de}\\ &X\longrightarrow Y\quad(\text{i.e. }\deleted{\text{les}}\text{ fibres de }\overline X\xrightarrow{\overline f}Y)\\ &\text{sont }n\text{-connexes, i.e. 0-connexes}\\ &\text{et }\pi_i(Z_y,z)=0\text{ pour }1\leq i\leq n.\end{aligned}\right.\] Alors $X$ \inserted{dans $\mathrm{Cat}$} sera $W'_n$-asphérique \unsure{sss|ssi} il est $W_{\deleted{\unsure{n}}\unsure{n|\infty}}$-asphérique, i.e. $n$-connexe. Donc par \inserted{la façon 2 de la} définition de $W'_n$, l'axiome $W(4\,\mathrm b'')$ est satisfait. Il faut croire que ça pèche ailleurs ! Peut-être déjà la saturation \[ \begin{aligned}&\text{Si }gf\text{ et }g\text{ sont dans }W'_i,\text{ il n'en est pas forcément}\\ &\text{de }\unsure{=|\text{même}}\text{ de }f,\end{aligned}\] car si $\bar g\bar f$ (et \deleted{$g$} donc $\bar g$) est surjectif, il n'en est pas forcément de même pour $\bar f$ ! Donc $W'_n$ ne mérite pas le nom de ``localiseur''.

Par contre, $W_n$ me paraît excellent, si on ne lui demande pas l'impossible -- de satisfaire à un ``axiome des fibrations'' qui semble bien \inserted{(à peu de choses près)} caractériser le seul $W_\infty$ parmi les ``bons'' localiseurs. En tous cas l'axiome de connexité $W(5)$ est satisfait, et l'axiome des limites filtrantes $W(6)$ aussi. Par ailleurs $W(1)$ \deleted{\illegible{1}} (et même la saturation forte), $W(2)$, $W(3)$\deleted{\unsure{,}} sont satisfaits par \deleted{\unsure{tout}} $W_{\mathbb D}$ associé à un dérivateur -- ce qui confirme à nouveau leur importance. Les axiomes $W(4)$ (l'axiome ``draconien'' des fibrations, $\pm$ spécial à $W_\infty$), $W(5)$ (connexité) et

% archive sheet number: 68
\scanpage{69}
\mspage{18}

$W(6)$ (limites) m'apparaissent comme de nature plus spéciale, je ne songerais à imposer aucun d'eux pour l'\deleted{\unsure{introduction}} \inserted{usage} terminologique de ``localiseur fondamental''.

II) Perplexités sur les axiomes et la terminologie

\deleted{A propos de la notion de $W$-fibration, quand $W$ provient d'un dérivateur, on a présumé que \deleted{\unsure{les morphismes}} \inserted{une} $W$-fibration gauche, ou droite, est déjà une $W$-fibration}

Dans I et la plupart du temps, en travaillant avec \unsure{un dérivateur|des dérivateurs}, j'ai supposé $\boxed{W_{\mathbb D}\supset W_\infty}$, \deleted{\unsure{donc}} hypothèse qui m'a \unsure{parue|paru} essentielle (à tort ou à raison).
% slanted
\margin{à \uline{tort}, il s'avère \unsure{maintenant}.} On voit donc que si $W_{\mathbb D}$ satisfait à l'axiome des fibrations\inserted{$W(4)$}, même sous sa forme ``la plus faible $W(4\,\unsure{\mathrm d|\mathrm b}'')$, il est \uline{égal} à $W_\infty$. Cela montre \deleted{\unsure{que}} \inserted{bien} que cet axiome doit être vu comme \deleted{\unsure{prati}}quement une \uline{caractérisation} de $W_\infty$ (parmi les localiseurs fondamentaux qui contiennent $W_\infty$ -- et ces localiseurs me paraissent les plus importants, et peut-être les seuls vraiment utiles.)

La théorie des $\mathbb D$-fibrations de I, p. 111--\unsure{126|125}, a été faite dans l'esprit qu'on vient de dire, en supposant $W_{\mathbb D}\supset W_\infty$. \deleted{\unsure{Cette notion}} \deleted{\unsure{avait été}} \inserted{La notion de $\mathbb D$-fibration} \inserted{à laquelle \unsure{je m'étais finalement arrêté|j'étais finalement arrivé} était} par suite différente des notions envisagées \deleted{\unsure{de}} ici dans un esprit différent (\unsure{profa.|prop.} 4). Elle revient à mettre la condition a) (trop forte presque tjrs, \deleted{\unsure{qu}} sauf pour $W=W_\infty$), caractérisant les $W$-fibrations ultrafortes, par des $Y''\longrightarrow Y'$ où $Y''$, $Y'$ sont tous deux $W_\infty$-asphériques (ce qui est plus \uline{\unsure{exigent|exigeant}} que $W$-asphériques) ; \deleted{\unsure{c'est seulement}} qd $Y''$ et $Y'$ sont $W_\infty$-asphériques \inserted{(au lieu de $W$-asphériques)} qu'on demande que \deleted{la condition que} $X''\longrightarrow X'$ soit \deleted{\unsure{une}} dans $W$. On trouve dans

% archive sheet number: 69
\scanpage{70}
\mspage{19}

une notion moins exigeante que celle de $W$-fibration forte. Si on dénote ici cette condition par b), on trouve (sous réserve que $W_\infty\subset W$), avec les notations de la p. 4

\diagram{\begin{tikzcd}[ampersand replacement=\&,column sep=small,row sep=small] \& \& \& {\vphantom b} \\ \& {a\quad W\text{-fibration ultraforte}}\arrow[d,Rightarrow] \\ \& {a'\quad W\text{-fibration forte}}\arrow[d,Rightarrow] \\ \& {b\quad W\text{-fibration}}\arrow[dl,Rightarrow]\arrow[dr,Rightarrow]\arrow[uuurr,bend right=65] \\ {a_1\quad W\text{-fibration gauche}} \& \& {a_2\quad W\text{-fibration droite}} \end{tikzcd}} \margin{Cf la bonne notion de $W$-fibration dans un contexte où on ne suppose pas $W_\infty\subset W$, dans \unsure{Prop.} et déf. 3 p. 51, 52}

\deleted{\unsure{De plus}, si} $a_1$) et $a_2$) étaient aussi appelées $\mathbb D$-fibrations cohomologiques resp. $\mathbb D$-fibrations homologiques.$^{*}$
% slanted
\margin{Non, c'est une notion encore plus forte, à \unsure{coups|coup} de car. \unsure{coexacts|exacts}, cf I p. 121.} Quand $\mathbb D$ satisfait Der 3 a$'$ (I p. 111) sur les $\mathbb D$-coefficients loc. \unsure{constants} sur les $X$ $W_\infty$-asphériques, alors les notions b, $a_1$, $a_2$ sont équivalentes. (Je présume que c'est le cas pour tous les $\mathbb D$ et tous les $W$ vraiment intéressants.) Par contre, j'ai l'impression que la notion de $W$-fibration forte, voire ultraforte, n'est guère utilisable que dans le cas $W=W_\infty$, où elle coïncide avec les précédentes. C'est dire qu'on peut probablement oublier la terminologie, et se rappeler par contre que dans le cas $W=W_\infty$, les $W$-fibrations ont des propriétés d'une force extrême, s'exprimant par les deux propriétés dégagées dans $W(4\,\mathrm a)$ et $W(4\,\mathrm b)$.

% archive sheet number: 70
\scanpage{71}
\mspage{20}

Dans la théorie de $\mathrm{HOT}_W$, influencé par la réflexion générale sur les dérivateurs et l'axiome Der 3 b, j'ai travaillé sous l'hypothèse $W\supset W_\infty$, mais en utilisant un axiome de fibration $\mathrm L\,5=W(4\,\mathrm b)$ qui, il s'avère, implique fatalement $W=W_\infty$. J'ai l'impression que je n'ai jusqu'à présent (propriétés de $f_!$ pour $\mathrm{HOT}$) pas utilisé $\unsure{V|W}\supset W_\infty$, donc ça peut être utile, comme ``ça ne coûte pas plus cher'' (et oblige d'être plus attentif à ce qu'on utilise), de laisser tomber $W\supset W_\infty$ \deleted{\unsure{et}} \deleted{\unsure{de de}} (impliquant $W=W_\infty$), et de mettre par contre sur $W$ des axiomes aussi forts (concernant \deleted{\unsure{sa}} fibration) que les propriétés de $W_\infty$, ``sans se priver en rien''. Je supposerai donc $W(1)$, $W(2)$, $W(3)$, $W(4\,\mathrm a)$ (la forme la plus forte de l'axiome des filtrations (cf pages 13, 14) -- c'est tout ce qui a été \unsure{utile|utilisé} jusqu'à présent. \deleted{Je} Si on y joint $W(5)$ (axiome de connexité) ça implique $W\subset W_\infty$. Mais pour l'instant je ne \unsure{crois|vois} pas en avoir un besoin. \deleted{Mais}

Mais méditant sur la question, je crois avoir trouvé le
% slanted
\margin{Cf pour un résultat plus fort XV \deleted{\illegible{2}} 9.12 p. \unsure{133bis3}.} \uline{Théorème}? (!) : Soit $W\subset\mathrm{Fl}(\mathrm{Cat})$ satisfaisant $W(1\deleted{\unsure{\mathrm b}},2,3,4\,\mathrm b'')$. Si $W$ ne satisfait pas à $W(5)$, alors $W=\mathrm{Fl}(\mathrm{Cat})$.
% slanted
\margin{\uline{NB} Le résultat \unsure{l.c.|cité} est plus général mais pas vraiment plus fort, car il ne donne pas $W=\mathrm{Fl}(\mathrm{Cat})$} \uline{Dém} : Soit $f:X\longrightarrow Y$ de $\mathrm{Cat}$, tel que $\pi_0(f)$ ne soit pas surjectif. Soit $I$ la catégorie discrète ayant pour \inserted{ens. d'}objets l'ens $\pi_0(Y)$, on a donc un diagramme commutatif \diagram{\begin{tikzcd}[ampersand replacement=\&,column sep=small,row sep=small] X\arrow[rr,"f"]\arrow[dr,"p"']\&\&Y\arrow[dl,"q"]\\ \& I\& \end{tikzcd}} d'autre part \inserted{un foncteur de} une catégorie \deleted{dans} une catégorie discrète est toujours $W_\infty$-parfait, donc $p,q$ sont des $W$-fibrations. Comme

% archive sheet number: 71
\scanpage{72}
\mspage{21}

% slanted
\margin{$\times$ on utilise \textcircled{4} sous forme plus forte que $W(4\,\mathrm b'')$ !} $f\in W$, il s'ensuit $f_i\in W$ pour tt $i\in I$.

Si $\pi_0(f)$ n'était pas surjective, prenant $i$ qui n'est pas en l'image, on aurait $X_i=\varnothing\longrightarrow Y_i\ne\varnothing$, qui serait dans $W$. Soit donc \[f:X=\varnothing\longrightarrow Y\] dans $W$. Alors $\forall Z$ dans $\mathrm{Cat}$, $f\times\mathrm{id}_Z:\underset{\substack{\Vert\\\varnothing}}{\underline{X\times Z}}\longrightarrow Y\times Z$ est dans $W$ (car $Z$ est $W$-fibrée sur $e$) donc $\varnothing\longrightarrow Y\times Z$ est dans $W$. Donc aussi la projection \diagram{\begin{tikzcd}[ampersand replacement=\&,column sep=small,row sep=small] Y\times Z\arrow[rr]\&\&Y,\\ \&\varnothing\arrow[ul]\arrow[ur]\&\end{tikzcd}} Comme cette projection est une $W$-fibration, il s'ensuit \deleted{que} par $W(4\,\mathrm b'')$ que ses fibres sont $W$-asphériques. Comme $Y\ne\varnothing$, il y a une telle fibre, elle est $\simeq Z$. Donc $Z$ est $W$-asphérique. \deleted{Soit $f:X$} Donc on trouve que \uline{toute} catégorie \inserted{$Z$ de $\mathrm{Cat}$} est $W$-asphérique. Donc toute flèche $Z\longrightarrow Z'$ de $\mathrm{Cat}$ est dans $W$, \unsure{cqed|cqfd}.

Supposons maintenant $\pi_0(f)$ surjective, \deleted{mais} non injective. Il y aurait une fibre de $\pi_0(X)\longrightarrow\pi_0(Y)=I$ qui serait de cardinal $>1$. En prenant le $f_i$ correspondant, on trouve $X_i\longrightarrow Y_i$ avec $Y_i$ 0-connexe, et $\mathrm{card}\,\pi_0(X_i)\geq2$. \deleted{\unsure{et}} Donc on est ramené au cas \[f:X\longrightarrow Y\qquad\left\{\begin{aligned}&f\in W\\&Y\text{ 0-connexe}\\&\mathrm{Card}\,\pi_0(X)\geq2.\end{aligned}\right.\]

Considérons la factorisation canonique \[X\xrightarrow{i_f}\overline X\xrightarrow{\overline f}Y,\] avec $i_f\in W$, donc $\overline f\in W$. Donc par $(W\,4\,\mathrm b'')$ les fibres de $\overline f$ sont $W$-asphériques. Mais on voit tout de suite, \deleted{que} en décomposant $\overline X$ suivant les comp. conn. $X_i$ de $X$, que ces fibres ont un $\pi_0$ de card $\geq\mathrm{Card}\,\pi_0(X)\geq2$. Donc on est ramené au cas où $Y=e$, i.e \[X\quad W\text{-asphérique,}\qquad \mathrm{card}\,\pi_0(X)\geq2.\] Considérons $\unsure{X}\times X$, qui est aussi $W$ asphérique, et l'application diagonale

% archive sheet number: 72
\scanpage{73}
\mspage{22}

\[d:X\longrightarrow X\times X=Y\] qui est donc dans $W$. \deleted{\unsure{Donc}} Mais \deleted{\unsure{alors}} $\pi_0(d)$ s'identifie à l'application diagonale \deleted{\unsure{pour}} $\pi_0(X)\longrightarrow\pi_0(X)\times\pi_0(X)$, qui n'est pas surjective puisque $\mathrm{card}\,\pi_0(X)\geq1$. Donc par la première partie, $W=\mathrm{Fl}(\mathrm{Cat})$.

Donc sous forme plus positive :

\uline{Corollaire} Soit $W\subset\mathrm{Fl}(\mathrm{Cat})$ qui satisfait à $W(1,2,3,4\,\mathrm b'')$, et telle que $W\ne\mathrm{Fl}(\mathrm{Cat})$. Alors $W$ \unsure{satisfais|satisfait} l'axiome de connexité $W(5)$, et par suite (th. 1, page 12) $W\subset W_\infty$.

III \unsure{Autres remarques|Autre remarque}. \unsure{Nom} des foncteurs liés à $W$.

Finalement, j'appellerai ``\uline{localiseur fondamental}'' une partie $W$ de $\mathrm{Fl}(\mathrm{Cat})$ satisfaisant les trois conditions suivantes :

(i) Saturation $W(1)$ (cf rappel p. 1)

(ii) Axiome de l'élément final $W(3)$ (page 3)

(iii) La variante affaiblie \inserted{$W(2\,\mathrm{bis})$} de \deleted{$W(2$} l'axiome de localisation $W(2)$ donnée dans le Cor. page 1 (anciennement c'était $\mathrm L\,4'$ -- dans P.S.)

Quand j'aurai besoin de plus sur $W$, je le dirai explicitement. Ces trois axiomes ont suffi pour \deleted{\unsure{que}} pratiquement tout ce que j'ai fait avec $W$ dans P.S., la seule exception étant un usage (si je me rappelle) unique, quelque part, d'une variante affaiblie de $W(2)$ (pour les foncteurs \uline{cartésiens} entre $S$-catégories \uline{fibrées} -- si \deleted{il induit} \unsure{dans} \unsure{toute} fibre il est $\in W$, il est lui-$=$ $\in W$).

On supposera donc par la suite donné un localiseur fondamental \deleted{$\unsure{W}$} \[W\subset\mathrm{Fl}(\mathrm{Cat}).\]

\uline{Proposition 1} Soit $f:X\longrightarrow Y$ \deleted{dans} \unsure{de} $\mathrm{Cat}$. Sont équiv. :

a) Pour tout \unsure{isom|isom.} local $Y'\longrightarrow Y$, $f':X'\longrightarrow Y'$ est $\in W$.

b) Même condition, dans le cas particulier $Y'=Y_{/y}\longrightarrow Y$, donc \deleted{$\unsure{\forall}$} $y\in Y$, $\underset{\unsure{\simeq}\,X\times_Y(Y_{/y})}{X_{/y}}$ est $W$-asphérique.

% archive sheet number: 73
\scanpage{74}
\mspage{23}

\uline{Définition 1} Sous les conditions \unsure{équ.|équiv.} a) b) de la prop 1, on dit que $f$ est \uline{$W$-asphérique}. On dit $f$ est \uline{$W$-coasphérique} si $f^\circ$ est $W$-asphérique, i.e. si les $ {}_{y}\backslash X$ sont $W$-asphériques (pour $y\in\mathrm{Ob}\,Y$).

\deleted{La condition \unsure{a)}} \inserted{\uline{Dém}} Il suffit de prouver $\mathrm b\Rightarrow\mathrm a$. Par définition d'un isom local, on voit que la condition b) est stable par changement de base $Y'\longrightarrow Y$, avec $Y'$ un isom local (puisque les $Y'_{/y'}\xrightarrow{\sim}Y_{/y}$, donc $X'_{/y'}\simeq X_{/y}$). Il suffit donc de voir que si $f$ \unsure{satisfait} à b), alors $f\in W$. Mais c'est l'axiome $W(2\,\mathrm{bis})$.

\deleted{\diagram{\begin{tikzcd}[ampersand replacement=\&,column sep=tiny,row sep=tiny] \&\&\deleted{X}\arrow[d]\\ \deleted{\unsure{Y_{/y'}=Y'}}\arrow[r]\&\deleted{\unsure{Y}}\&\deleted{Y}\&\deleted{Y'}\arrow[l]\\ \deleted{Y'_{/y'}}\arrow[u]\arrow[r]\&\deleted{Y_{/y}}\arrow[u]\&\deleted{Y_{/y}}\arrow[u]\&\deleted{Y'_{/y'}}\arrow[l]\arrow[u] \end{tikzcd}}} \diagram{\begin{tikzcd}[ampersand replacement=\&,column sep=small,row sep=small] X\arrow[d]\&\&\\ Y\&Y'\arrow[l]\&\\ \&Y_{/y}\arrow[ul]\&Y'_{/y'}\arrow[ul]\arrow[l] \end{tikzcd}}

\uline{Remarque} Si $f$ a un adjoint à droite $g$, alors pour tout $y\in Y$, $X_{/y}\simeq X_{/g(y)}$, donc $X_{/y}$ a un objet final et est donc $W$-asphérique. Donc si $f$ a un adjoint à droite (à gauche), alors $f$ est $W$-asphérique ($W$-coasphérique).

Mais $=$ si $f$ n'a pas d'adjoint à droite, il y a \unsure{cepdt} $=$ un adjoint généralisé \[g:Y\longrightarrow X^\wedge,\] et on a \[X_{/y}\simeq X_{/g(y)}\qquad(\text{pour l'inclusion tt. fid. }X\deleted{\unsure{{}^{\wedge}}}\hookrightarrow X^\wedge),\] donc $\mathrm{Top}(X_{/y})$ s'identifie au topos induit $\mathrm{Top}(X)_{/g(y)}$. La condition $f$ $W$-asphérique signifie donc que les \inserted{faisceaux} $g(y)$ \deleted{dans} $\mathrm{Top}(X)$ sont

% archive sheet number: 74
\scanpage{75}
\mspage{24}

$W$-asphériques. Pour de nombreuses \deleted{autres} \deleted{\unsure{équivalentes}} formulations équivalentes, je renvoie à P.S.

\uline{Proposition 2} Soit $X\mathrel{\overset{f}{\relbar}}Y\xrightarrow{g}Z$ dans $\mathrm{Cat}$. Si $f$ et $g$ sont \deleted{\unsure{des}} $W$-asphériques (\inserted{$W$-}coasphériques), de $=$ $gf$. Si $gf$ et $f$ sont $W$-asphériques ($W$-coasphériques), $g$ aussi.

\deleted{Rappel} \uline{Rappel} Un morphisme $p:S'\longrightarrow S$ de $\mathrm{Cat}$ est un \uline{isom. local} si $\forall\,s'\in\mathrm{Ob}\,S'$, le morphisme induit \[p_{s'}:S'_{/s'}\xrightarrow{\sim}S_{/s}\qquad(s=p(s'))\] est un isomorphisme : $p$ induit un isomorphisme entre les ``localisées'' de $S'$ et $S$ en $s'$ et $s=p(s')$ respectivement. Il est immédiat que cette notion est stable par composition, que si $gf$ et $g$ sont des isom locaux, alors $f$ l'est aussi, que tout isomorphisme est un isom. local. De plus, les isom. locaux sont \uline{stables par chgt de base}, ce qui résulte du fait \inserted{(trivial et important)} que si on a un carré cartésien de $\mathrm{Cat}$ \diagram{\begin{tikzcd}[ampersand replacement=\&] X\arrow[d,"f"']\& X'\arrow[l,"q"']\arrow[d,dash,"f'"]\\ S\& S'\arrow[l,"p"] \end{tikzcd}}, et si $x'\in\mathrm{Ob}\,X'$ a pour images $x,s',s$ dans $X,S',S$, alors le diagramme \diagram{\begin{tikzcd}[ampersand replacement=\&] X_{/x}\arrow[d,"{f_{/x}}"']\& X'_{/x'}\arrow[l,"{q_{/x'}}"']\arrow[d,"{f'_{/x'}}"]\\ S_{/s}\& S'_{/s'}\arrow[l,"{p_{/s'}}"] \end{tikzcd}} est également cartésien. Donc si $p_{/s'}$ est iso, $q_{/x'}$ aussi.

% archive sheet number: 75
\scanpage{76}
\mspage{25}

Revenons à la prop. 2. La stabilité de la notion de morphisme asphérique par composition \diagram{\begin{tikzcd}[ampersand replacement=\&,row sep=small] X\arrow[d,"f"']\&X'\arrow[l,"r"']\arrow[d,"f'"]\\Y\arrow[d,"g"']\&Y'\arrow[l,"q"']\arrow[d,"{g'}"]\\ Z\&Z'\arrow[l,"p"]\end{tikzcd}} résulte aussitôt de la forme a) de la condition. De même pour la condition que si $gf$ et $f$ sont $W$-asphériques, $g$ aussi -- il faut montrer que si $Z'\longrightarrow Z$ est iso loc., alors $g'\deleted{\unsure{f'}}$ est dans $W$, or l'hyp. implique que $g'f'$ et $f'$ le sont.

\deleted{\uline{NB} \unsure{Quand}} \uline{Corollaire 1} L'image inverse d'un $f:X\longrightarrow Y$ qui est $W$-asphérique \inserted{(resp. $W$-coasphérique)} par un iso local (resp. par un co-iso local) est $W$-asphérique ($W$-coasphérique)

\diagram{\begin{tikzcd}[ampersand replacement=\&,column sep=small] X\arrow[d,dash]\&X'\arrow[l,dash]\arrow[d,dash]\&X''\arrow[l]\arrow[d]\\ \unsure{Y}\&Y'\arrow[l,"{\mathrm{iso\ loc}}"]\&Y''\arrow[l,"{\mathrm{iso\ loc}}"]\end{tikzcd}} Résulte formellement de la définition par a), et du fait qu'un composé de deux iso. locaux (resp coisom locaux) est itou.

\uline{NB} Bien sûr, la notion de foncteur $W$-asphérique n'est pas stable par chgt. de base quelconque. \deleted{Mais} Par exemple, l'inclusion $X=\{y_0\}\mathrel{\relbar}Y$ est $W$-asph. \unsure{sss|ssi} les catégories discrètes $\mathrm{Hom}\,(y_0,y)$ sont $W$-asph. i.e. (sous réserve \unsure{que l'on ait|qu'elles ont} $W(5)$, axiome de connexité) sss les $\mathrm{Hom}(y_0,y)$ sont de cardinal 1, i.e. $y_0$ \uline{objet initial} de $Y$. Mais si $Y\ne\{y_0\}$, prenant $y$ dans $Y$, $y\ne y_0$, \deleted{la fibre de} et prenant $Y'=\{y\}\longrightarrow Y$, \unsure{base|bas.} $X'=X\times_Y Y'=\varnothing$, et $\varnothing\longrightarrow e$ n'est pas asphérique !

% archive sheet number: 76
\scanpage{77}
\mspage{26}

Par contre, on s'attend que les \deleted{\unsure{notions}} morphismes $W$-asphériques soient stables par chgt. de base $W$-\uline{lisses} (qui généralisent les iso. locaux), et les morphismes $W$-coasphériques par chg. de base $W$-\unsure{propres|propre}. Cf plus bas.

\uline{Corollaire 2} Soit $f:X\longrightarrow Y$ \inserted{dans $\mathrm{Cat}$} une \uline{équivalence} de catégories. Alors $f$ est $W$-asphérique et $W$-coasphérique.

En effet, $f$ a un adjoint à \unsure{g} et un adjoint à droite, et on utilise la remarque page 23.

\uline{Prop. 3}. Soit $f:X\longrightarrow Y$ de $\mathrm{Cat}$. Conditions équivalentes :

a) Pour tout $y_0\in\mathrm{Ob}\,Y$, le foncteur d'inclusion \[X_{y_0}\mathrel{\relbar}X_{/y_0}\tag*{(\unsure{*|x})}\] est $W$-\uline{co}asphérique.

b) Pour tout $x_0\in X$, le morphisme induit \unsure{par les colocalisés|pour les colocalisés} \[{}_{x_0}\backslash X\mathrel{\relbar}{}_{y_0}\backslash Y\qquad(\text{où }y_0=f(x_0))\] est à fibres $W$-asphériques.

c) Pour tt \deleted{$\unsure{\Delta}$} morphisme $\mathord{\tikz[baseline=0pt,x=1ex,y=1ex]{\draw[line width=.55pt,line join=round] (0,0)--(.78,1.4)--(1.53,0)--cycle; \draw[line width=.55pt] (.25,0)--(.90,1.18);}}^1=Y'\longrightarrow Y$, considérant $X'=X\times_Y Y'$ sur $\mathord{\tikz[baseline=0pt,x=1ex,y=1ex]{\draw[line width=.55pt,line join=round] (0,0)--(.78,1.4)--(1.53,0)--cycle; \draw[line width=.55pt] (.25,0)--(.90,1.18);}}^1$, l'inclusion $X'_1\longrightarrow X'$ de la fibre fermée de $X'$ (fibre sur le point fermé de $Y'$) est $W$-coasphérique.

d) Pour toute flèche $u_0:y_0\longrightarrow y_1$ dans $Y$, et tout $x_0$ de $X_{y_0}$, la catégorie $\unsure{X|\mathscr X}(x_0,u_0)$ formée des flèches \inserted{$u:x_0\to x$} dans $X$ de source $x_0$, et qui relèvent $u_{\unsure{0}}$, est $W$-asphérique.

% archive sheet number: 77
\scanpage{78}
\mspage{27}

\uline{Dém}. Pour expliciter la condition a), il faut \deleted{\unsure{o}} \deleted{expliciter les fibres de \unsure{ce mor}} expliciter la catégorie $  {}_{\unsure{\xi|\zeta}_0}\backslash(X_y)$, pour un objet quelconque \inserted{$\xi_0=$} $(x_0,u_0)$ de $X_{/y}$, \deleted{i.e. que une flèche} \deleted{quelconque de but $y$ de $X$} \inserted{où} $x_0\in\mathrm{Ob}\,X$, \deleted{\unsure{et}} \[u_0:y_0\longrightarrow y_1=y\quad\text{de }Y,\text{ avec }y_0=f(x_0).\] Cette catégorie $  {}_{\xi_0}\backslash X_{y_1}$ s'explicite alors comme la catégorie des $x_1$ dans $X_{y_1}$, \uline{munis} d'une flèche \inserted{de $\xi_0$} dans \deleted{\unsure{la}} l'image \inserted{$\xi_1=(x_1,\mathrm{id}_{f(x_1)}\unsure{=|\Rightarrow}\mathrm{id}_{y_1})$} de $x_1$ dans $X_{/y_1}$ de $\xi_0$, i.e. d'une flèche $u$ de $x_0$ dans $x_1$, \diagram{\begin{tikzcd}[ampersand replacement=\&,column sep=small,row sep=small] x_1\arrow[d,dash]\&\&x_0\arrow[ll,"u"']\\ y_1\&y_1\arrow[l,dash,"{\mathrm{id}}"']\&\\ \&\&y_0\arrow[uu]\arrow[ul,"{u_0}"']\arrow[ull,"{f(u)}"] \end{tikzcd}} telle que $f(u)$ donne lieu au triangle comm. dans $Y$, i.e. $f(u)=u_0$. Donc c'est la catégorie notée $X(x_0,u_0)$ dans d). On voit tout de suite que \deleted{c'est} \inserted{elle est} aussi \inserted{iso à} la fibre de $  {}_{x_0}\backslash X\mathrel{\relbar}{}_{y_0}\backslash Y$ en l'élément $u_0$ de $ {y_0}\backslash Y$, dont il est question dans b). D'où l'équivalence de a) b) \unsure{c)|d)}.

D'autre part, la donnée de $u_0$ dans $Y$ équivaut à celle d'un foncteur $\mathord{\tikz[baseline=0pt,x=1ex,y=1ex]{\draw[line width=.55pt,line join=round] (0,0)--(.78,1.4)--(1.53,0)--cycle; \draw[line width=.55pt] (.25,0)--(.90,1.18);}}^1\mathrel{\relbar}Y$, \unsure{comme} dans c), et pour exprimer alors que $\underset{\substack{\mathord{\tikz[baseline=0pt,x=1ex,y=1ex]{\draw[line width=.45pt] (0,0) .. controls (-.28,.45) and (.28,.65) .. (0,1); \draw[line width=.45pt] (.30,0)--(.30,1);}}\\X_{y_1}}}{X'_1}\longrightarrow X'$ est coasphérique, il suffit de \deleted{\unsure{vérifier}} \inserted{l'exprimer} que les $  {}_{x'}\backslash X'_1$ \inserted{($x'\in\mathrm{Ob}\,X'$)} sont $W$-asphérique. Or c'est trivial si $x'\in\mathrm{Ob}\,X'_1$ (objet initial), donc il suffit de le vérifier pour $x'$ \deleted{$\in$} \inserted{dans} $X'\smallsetminus X'_1=X'_0\simeq X_{y_0}$. \deleted{\unsure{Et}} \inserted{Alors $x'$ s'identifie à un élément $x_0$ de $X_{y_0}$,} on vérifie tout de suite que $  {}_{x_0}\backslash X'_1$ est isom à la catégorie $X(x_0,u_0)$ de d). \unsure{Céla|Cela} achève la dém.

% archive sheet number: 78
\scanpage{79}
\mspage{28}

\uline{Définition 2} Sous les conditions de la prop. 3, on dit que $f$ est $W$-\uline{propre}. On dit que $f$ est $W$-lisse si $f^\circ$ est $W$-propre, i.e. si les morphismes d'\unsure{inclusions|inclusion} (\unsure{*|x}) sont \uline{$W$-asphériques} etc (conditions duales de celles de la prop 3).

\uline{Remarque} La condition a) de la prop 3 est satisfaite en particulier si l'inclusion \inserted{(*)} a un adjoint à gauche (cf remarque page 23). Cela signifie que les catégories $X(x_0,u_0)$ ont un \unsure{objet} initial. Cela signifie aussi que la catégorie $X$ sur $Y$ est $\mathrm{Cat}$-cofibrée sur $Y$ au sens faible, i.e. qu'il existe, pour $u_0:y_0\mathrel{\relbar}y_1$ flèche dans $Y$, \deleted{un foncteur} et \unsure{$a_0$|$x_0$} dans $X_{y_0}$, un objet $u_{0*}(x_0)$ dans $X_{y_1}$ qui représente le foncteur covariant $X_{y_1}\mathrel{\relbar}(\mathrm{Ens})$, $x_1\longrightarrow\mathrm{Hom}_{\unsure{u|u_0}}(x_0,x_1)$. Ainsi, si $X$ est $\mathrm{Cat}$-cofibrée au sens faible sur $Y$, et a fortiori si elle est $\mathrm{Cat}$-cofibrée, alors $f:X\longrightarrow Y$ est $W$-propre. Dualement, si $f$ est faiblement $\mathrm{Cat}$-fibrée, $f$ est lisse.

% slanted
\margin{\unsure{Dire aussi|D'où, aussi|D'où, on voit} : Si $X\xrightarrow{f}Y\xrightarrow{g}Z$, si $gf$ et $f$ $W$-lisses, alors $g$ est lisse (en les points de $\mathrm{Im}\,f$ ?), idem pour propre.} \uline{Proposition 4} \deleted{Si} La notion de morphisme \deleted{\unsure{fa}}propre est stable par composition, par changement de base. Les coisom locaux sont propres, \inserted{(\unsure{trivial dans b|trivial selon b})} a fortiori les isom, sont propres, les imm. fermées sont propres.

Enoncé dual pour les morphismes \unsure{lisse|lisses} : \unsure{stabilité} par chgt de base, par composition, les iso. locaux sont lisses, notamment les imm. ouvertes.

% archive sheet number: 79
\scanpage{80}
\mspage{29}

La stabilité par chgt. de base est immédiate sur la forme c). \deleted{\unsure{La}} Stabilité par composition moins triviale : \diagram{\begin{tikzcd}[ampersand replacement=\&,row sep=small,column sep=small] {x\quad X}\arrow[d,"f"]\&{{}_{x}\backslash X}\arrow[l,Rightarrow]\arrow[d]\&\overline X\arrow[l]\arrow[d]\\ {y\quad Y}\arrow[d,"g"]\&{{}_{y}\backslash Y}\arrow[l]\arrow[d]\&{\deleted{\unsure{Z}}\overline Y}\arrow[l]\arrow[d]\\ {z\quad Z}\&{{}_{z}\backslash Z}\arrow[l]\&\{u\}\arrow[l] \end{tikzcd}} Dans le diagramme à carrés cart. ci-contre (déduit de $x\in\mathrm{Ob}\,X$, d'où $y=f(x)$, $z=g(y)=gf(x)$\deleted{\unsure{)}}, et de $u\in\mathrm{Ob}\,{}_{z}\backslash Z$), il faut prouver que $\overline X$ est $W$-asphérique. On sait que $\overline Y$ l'est ($g$ propre) et que $\overline X\mathrel{\relbar}\overline Y$ est à fibres $W$-asphériques. Il suffit donc d'utiliser la prop. 5 ci-dessous, et la remarque immédiate :

\uline{Lemme} Si $f$ est propre, il en est de $=$ de \deleted{qui est} $ {}_{x}\backslash f:{}_{x}\backslash X\longrightarrow {}_{y}\backslash Y$, \inserted{pour tt $x$ de $X$, $y=f(x)$,} En effet, si $u'$ est dans $ {}_{x}\backslash X$ et $v'$ dans $ {}_{y}\backslash Y$, \deleted{alors} et si $x',y'$ sont leurs images sur $X$, $Y$ resp., alors \[{}_{u'}\backslash({}_{x}\backslash X)\simeq {}_{x'}\backslash X,\qquad {}_{v'}\backslash({}_{y}\backslash Y)\simeq {}_{y'}\backslash Y\quad\deleted{\text{et}}\] et la flèche entre les deux s'identifie à la flèche de localisation \[{}_{x'}\backslash X\longrightarrow {}_{y'}\backslash Y,\]

% slanted
\margin{\deleted{En fait, si $f$ propre, $f$ est asphérique, \unsure{si} $f$ lisse, $f$ est \unsure{coasph.}}} \uline{Prop 5} Soit $f:X\longrightarrow Y$ de $\mathrm{Cat}$, supposons $f$ $W$-propre \deleted{et} $W$-lisse, et à fibres $W$-asphériques, alors $f\in W$ (et $=$ $f$ est \uline{universellt} dans $W$, i.e. il est dans $W$ et le reste par \uline{tout} chgt. de base).

% archive sheet number: 80
\scanpage{81}
\mspage{30}

Comme les hyp. sont stables par chang. de base, il suffit de prouver $f\in W$. Supposons p. ex. $f$ propre, prouvons que $f$ est $W$-\uline{asphérique}, i.e. que les $X_{/y}$ ($y\in\mathrm{Ob}\,Y$) sont $W$-asphériques. Mais $X_y\longrightarrow X_{/y}$ étant $W$-coasphérique par hyp. de propreté, est $\in W$ (par $W(2\,\mathrm{bis})$) et comme $X_y$ est $W$-asph. par hyp., $X_{/y}$ aussi. Mais par $W(2\,\mathrm{bis})$, $f$ $W$-asphérique $\Longrightarrow f\in W$. Le cas $f$ lisse est dual.

\deleted{\uline{Remarque}} Ainsi, \deleted{\unsure{si}} $f$ est propre, s'il est à fibres $W$-asphériques il est ``universellement $W$-asphérique'', i.e. est et reste $W$-asphérique par \deleted{\unsure{tout}} ch. \deleted{\unsure{gement}} de base. Inversement, si $f$ est univ. \inserted{dans} $W$\deleted{-asphérique} (pas forcément propre) alors \inserted{\unsure{en particulier|en vertu de}} ses fibres sont $W$-asphériques \deleted{Cela résulte du fait que \deleted{\unsure{$\xrightarrow{f}$}} $Z\longrightarrow e$ est $W$-asphérique (ou $W$-coasphérique) \unsure{sss|ssi} $Z$ est $W$-asphérique (trivial). Notons que $f$ peut être universellement $W$-asphérique c'est \unsure{-- à la fois univ. $W$-asphér.}} Mais 5. n'implique bien sûr pas que $f$ soit propre, ou lisse.

% slanted
\margin{Mais ``universellt dans $W$'' implique déjà ``universellement $W$-asphérique'', \unsure{NB} la notion de $W$-asphéricité est définie par prop. 1 a).}

\uline{Corollaire} Si $W$ satisfait $W(2)$ (pas seulement $W(2\,\mathrm{bis})$), alors pour un $S$-morphisme $X\longrightarrow Y$ de cat. qui sont toutes deux \deleted{\unsure{$W$-}}propres sur $S$, ou \deleted{\unsure{toutes}} deux $W$-lisses, si $f_s:X_s\mathrel{\relbar}Y_s$ est dans $W$ pour tt $s\in S$, alors $f\in W$.

% archive sheet number: 81
